\section{Related Work}\label{sec:related}

Classical statistical methods (ARIMA, Holt-Winters) capture linear temporal dependencies but are consistently outperformed by machine learning when exogenous weather predictors matter: SVR and MARS beat ARIMA on short-term Australian demand \parencite{almusaylh2018forecasting}, and tree-based methods dominate linear regression for residential consumption \parencite{schirmer2019evaluation}. Among tree ensembles, Random Forest \parencite{breiman2001randomforest} reduces variance through bagging; XGBoost \parencite{chen2016xgboost} adds regularised sequential boosting and is widely used in energy applications \parencite{khan2020hybrid}; CatBoost \parencite{prokhorenkova2018catboost} contributes ordered boosting and native categorical-feature handling, which suits demand data rich in calendar categories. Weather, holidays, and weekends are recognised as key influencing factors \parencite{khan2021influencing}.

Meta-heuristic hyperparameter tuning has shown large gains---e.g.\ 42\% accuracy improvement from PSO-optimised models on university buildings \parencite{li2020datadrivenstrategy}. PPSO \parencite{ghasemi2019ppso} reformulates PSO with a phase angle that modulates cognitive and social terms through periodic functions, eliminating manual tuning of inertia and acceleration coefficients. Zhang et al.\ \parencite{zhang2023catboostppso} combined CatBoost with PPSO for hourly consumption of a single customer and found it superior to five alternative optimisers.

\textbf{Gap.} Prior work tunes only a \emph{single} favoured model, leaving open whether reported gains reflect the algorithm or preferential tuning, and the CatBoost-PPSO methodology has not been evaluated on daily national-level demand in a tropical climate. This project addresses both gaps.
